\documentclass[11pt,a4paper]{article}

% ------------------------------------------------
% Packages
% ------------------------------------------------

\usepackage[T1]{fontenc}
\usepackage[utf8]{inputenc}
\usepackage{lmodern}
\usepackage{microtype}

\usepackage[
    top=2.2cm,
    bottom=2.2cm,
    left=2.4cm,
    right=2.4cm
]{geometry}

\usepackage{xcolor}
\usepackage{enumitem}
\usepackage{titlesec}
\usepackage{fancyhdr}
\usepackage{lastpage}

% ------------------------------------------------
% Colors
% ------------------------------------------------

\definecolor{accent}{HTML}{244A73}
\definecolor{textgray}{HTML}{555555}
\definecolor{lightgray}{HTML}{D9DEE5}

% ------------------------------------------------
% General layout
% ------------------------------------------------

\setlength{\parindent}{0pt}
\setlength{\parskip}{7pt}

\setlist[itemize]{
    leftmargin=1.6em,
    itemsep=3pt,
    topsep=4pt,
    parsep=0pt
}

% ------------------------------------------------
% Section formatting
% ------------------------------------------------

\titleformat{\section}
    {\Large\bfseries\color{accent}}
    {}
    {0pt}
    {}
    [\vspace{3pt}\color{lightgray}\titlerule]

\titleformat{\subsection}
    {\large\bfseries}
    {}
    {0pt}
    {}

\titlespacing*{\section}
    {0pt}{22pt}{10pt}

\titlespacing*{\subsection}
    {0pt}{14pt}{4pt}

% ------------------------------------------------
% Header and footer
% ------------------------------------------------

\pagestyle{fancy}
\fancyhf{}

\fancyhead[L]{\small\color{textgray}[P05] Food Image Classification}
\fancyhead[R]{\small\color{textgray}Project Plan}

\fancyfoot[C]{
    \small\color{textgray}
    \thepage\ / \pageref{LastPage}
}

\renewcommand{\headrulewidth}{0.4pt}
\renewcommand{\footrulewidth}{0pt}

% ------------------------------------------------
% Document
% ------------------------------------------------

\begin{document}

% ------------------------------------------------
% Title page
% ------------------------------------------------

\begin{titlepage}

\thispagestyle{empty}

\vspace*{2.2cm}

{\color{accent}\rule{\textwidth}{1.5pt}}

\vspace{1.1cm}

{\Large\color{textgray}\textbf{[F26] Introduction to Computer Vision}}

\vspace{0.8cm}

{\Huge\textbf{Food Image Classification}}

\vspace{0.45cm}

{\LARGE\color{accent}\textbf{Project Plan}}

\vspace{1cm}

{\color{lightgray}\rule{\textwidth}{0.8pt}}

\vspace{1.4cm}

{\large
The project studies food image classification on a selected subset of Food-101.
The main experiments evaluate the effect of ImageNet pretraining and data augmentation
on classification quality and on errors between visually similar food classes.
}

\vfill

{\large\textbf{Dataset}}\\
Food-101

\vspace{0.5cm}

{\large\textbf{Main model}}\\
ResNet-18

\vspace{0.5cm}

{\large\textbf{Main experimental factors}}\\
ImageNet pretraining and data augmentation

\vfill

{\color{textgray}October 2026}

\end{titlepage}

\setcounter{page}{1}

% ------------------------------------------------
% Project Scope
% ------------------------------------------------

\section{Project Scope}

\subsection{Objective}

The goal of the project is to build and evaluate an image classifier for food recognition.
The project will use a subset of 12 classes from the Food-101 dataset:

\begin{itemize}
    \item carrot cake
    \item cheesecake
    \item chocolate cake
    \item red velvet cake
    \item ramen
    \item pho
    \item spaghetti bolognese
    \item spaghetti carbonara
    \item fried rice
    \item risotto
    \item hamburger
    \item pizza
\end{itemize}

\subsection{Class Selection}

Several visually similar classes are included on purpose. The selected subset contains
multiple cakes, noodle dishes, pasta dishes, and rice dishes.

This makes it possible to measure not only overall classification quality, but also
whether the model systematically confuses classes with similar appearance, ingredients,
texture, or presentation.

Pizza and hamburger are included as more visually distinctive classes. They provide
a useful comparison with the more difficult groups and help show whether the errors
are concentrated mainly among similar classes.

\subsection{Model and Experimental Factors}

The main model will be \textbf{ResNet-18}.

The project will study two factors:

\begin{itemize}
    \item ImageNet pretraining
    \item data augmentation
\end{itemize}

The same model architecture will be used in all main experiments. This keeps the
architecture constant and makes it possible to measure the effect of pretraining
and augmentation directly.

\subsection{Research Question}

\textbf{How do ImageNet pretraining and data augmentation affect classification
performance on a subset of Food-101 that contains visually similar food classes?}

% ------------------------------------------------
% Dataset Preparation
% ------------------------------------------------

\section{Dataset Preparation}

\subsection{Dataset}

The project will use the public Food-101 dataset. Only the 12 selected classes will
be included in the experiments.

Food-101 provides 750 training images and 250 test images for each class.

The original test set will remain unchanged. The original training set will be divided
into 600 training images and 150 validation images per class.

The final dataset will contain:

\begin{itemize}
    \item \textbf{7,200} training images
    \item \textbf{1,800} validation images
    \item \textbf{3,000} test images
    \item \textbf{12,000} images in total
\end{itemize}

The same split will be used in every experiment. This keeps the evaluated data constant,
so differences between experimental conditions are not caused by different samples.

\subsection{Image Preprocessing}

All images will be prepared for the 224 by 224 input size used by ResNet-18.

For experiments without augmentation, preprocessing will include:

\begin{itemize}
    \item resizing
    \item center cropping to 224 by 224 pixels
    \item conversion to tensors
    \item normalization
\end{itemize}

\subsection{Data Augmentation}

For experiments with augmentation, training images will use:

\begin{itemize}
    \item random resized cropping to 224 by 224 pixels
    \item random horizontal flipping
    \item small random color changes
    \item conversion to tensors
    \item normalization
\end{itemize}

These transformations create additional variation in viewpoint, framing, and color
while preserving the food class. This allows the model to see more varied versions
of the available training images.

Validation and test images will use fixed preprocessing so that evaluation does not
depend on random image transformations.

% ------------------------------------------------
% Experimental Plan
% ------------------------------------------------

\section{Experimental Plan}

\subsection{Experimental Conditions}

Four conditions will be compared:

\begin{itemize}
    \item ResNet-18 with random initialization and no augmentation
    \item ResNet-18 with ImageNet pretraining and no augmentation
    \item ResNet-18 with random initialization and augmentation
    \item ResNet-18 with ImageNet pretraining and augmentation
\end{itemize}

The first condition will serve as the baseline.

Comparing the first and second conditions will measure the effect of ImageNet
pretraining.

Comparing the first and third conditions will measure the effect of data augmentation.

The fourth condition will show how pretraining and augmentation perform when used
together.

\subsection{Training Strategy}

For pretrained models, the complete ResNet-18 network will be fine tuned.

This allows both the final classifier and the feature extraction layers to adapt from
general ImageNet features to features that are useful for food recognition.

The planned training settings are:

\begin{itemize}
    \item AdamW optimizer
    \item initial learning rate of $3 \times 10^{-4}$
    \item batch size of 32
    \item maximum of 30 epochs
    \item learning rate scheduling during training
    \item model selection based on validation accuracy
\end{itemize}

AdamW provides weight decay during optimization and can be used consistently for both
randomly initialized and pretrained models.

A batch size of 32 provides a practical balance between memory use and the number of
samples used for each gradient update.

The same training settings will be used across the four conditions so that the main
difference between experiments remains the use of pretraining and augmentation.

\subsection{Repeated Experiments}

Each condition will be trained three times with different random seeds.

Repeated runs reduce the influence of random initialization, training sample order,
and random augmentation on the final comparison.

The final results will therefore be reported using the \textbf{mean} and
\textbf{standard deviation} across the three runs.

\subsection{Evaluation Metrics}

The main evaluation metrics will be:

\begin{itemize}
    \item \textbf{Accuracy}. The selected classes contain the same number of images,
    so accuracy provides a direct measure of overall classification performance.

    \item \textbf{Macro F1 score}. The F1 score will be calculated separately for
    every class and then averaged. This helps identify whether strong overall
    performance hides weaker results for individual classes.
\end{itemize}

Additional analysis will include:

\begin{itemize}
    \item precision for each class
    \item recall for each class
    \item confusion matrix
    \item analysis of the most common confusion pairs
    \item manual inspection of representative classification errors
\end{itemize}

\subsection{Confusion and Failure Analysis}

The confusion matrix will be used to identify which classes are most frequently
confused.

This analysis is especially important for the selected groups of visually similar
foods, such as cakes, noodle dishes, pasta dishes, and rice dishes.

Around six representative errors will be inspected manually.

Possible causes will include:

\begin{itemize}
    \item similar ingredients
    \item similar shape or texture
    \item similar presentation
    \item unusual viewing angles
    \item background objects
    \item partial visibility of the food
    \item images that are visually ambiguous even for a human observer
\end{itemize}

The test set will be used only after training and model selection are complete.
Validation data will be used for checkpoint selection and training decisions.
The test set will provide the final evaluation.

% ------------------------------------------------
% Implementation Plan
% ------------------------------------------------

\section{Implementation Plan}

\subsection{Machine Learning Stack}

The main machine learning stack will consist of:

\begin{itemize}
    \item Python
    \item PyTorch
    \item torchvision
    \item scikit-learn
    \item pandas
    \item Matplotlib
\end{itemize}

PyTorch will be used for model training and inference.

Torchvision provides the Food-101 dataset interface, the ResNet-18 implementation,
ImageNet pretrained weights, and the image transformations required for preprocessing
and augmentation.

Scikit-learn will be used for classification metrics and confusion matrix calculation.

Pandas will be used to organize and compare experimental results.

Matplotlib will be used to create result visualizations.

\subsection{Model Development}

The implementation will include the following steps:

\begin{itemize}
    \item prepare the selected Food-101 classes
    \item create the fixed training and validation split
    \item implement preprocessing and augmentation
    \item configure ResNet-18 for 12 output classes
    \item train all four experimental conditions with three random seeds
    \item save model checkpoints and training results
    \item evaluate the trained models on the test set
    \item calculate the main metrics and metrics for each class
    \item generate the confusion matrix
    \item collect misclassified samples
    \item aggregate the results from repeated runs
    \item select the final model based on the experimental results
\end{itemize}

Training and evaluation will be implemented as Python programs. This format is suitable
for repeated machine learning experiments because the same training and evaluation
procedure can be executed with fixed configurations and random seeds.

\subsection{Prediction Backend}

The prediction backend will use \textbf{FastAPI}.

FastAPI allows the trained PyTorch model and the image preprocessing pipeline to be
used directly inside the Python backend while providing an HTTP interface for the
frontend.

The backend will:

\begin{itemize}
    \item load the final trained model when the application starts
    \item receive an image
    \item validate the input image
    \item apply the same preprocessing used during evaluation
    \item run model inference
    \item return the predicted class and confidence score
\end{itemize}

\subsection{Frontend}

The frontend will be implemented with \textbf{React}.

It will provide:

\begin{itemize}
    \item image upload
    \item image preview
    \item communication with the backend
    \item predicted food class
    \item prediction confidence
    \item loading states
    \item error handling
\end{itemize}

\subsection{Infrastructure}

The project infrastructure will use:

\begin{itemize}
    \item GitHub
    \item Docker
    \item GitHub Actions
\end{itemize}

Docker will provide a consistent environment for the application and its dependencies.

GitHub Actions will automatically check that the project can be installed, tested,
and built correctly after changes to the repository.

The experiment setup will also preserve the configurations and random seeds required
to reproduce the reported training and evaluation results.

% ------------------------------------------------
% Team Responsibilities
% ------------------------------------------------

\section{Team Responsibilities}

\subsection{Denis Geraskin}

\textbf{Email:} d.geraskin@innopolis.university

\textbf{Area: AI Engineering and Backend}

Responsibilities:

\begin{itemize}
    \item implementation and configuration of ResNet-18
    \item implementation of random and ImageNet pretrained initialization
    \item training and fine tuning of the model
    \item execution of the main experimental conditions
    \item model checkpoint handling
    \item implementation of model inference
    \item model related backend logic
\end{itemize}

\subsection{Denis Beliaev}

\textbf{Email:} d.beliaev@innopolis.university

\textbf{Area: Data Science, Data Preparation and Backend}

Responsibilities:

\begin{itemize}
    \item preparation of the Food-101 subset
    \item creation of the training, validation, and test split
    \item implementation of preprocessing and augmentation
    \item implementation of quantitative evaluation
    \item calculation of Accuracy, Macro F1, precision, and recall
    \item generation and analysis of the confusion matrix
    \item aggregation of results from repeated runs
    \item identification of important confusion patterns
    \item backend image preprocessing and validation
\end{itemize}

\subsection{Aiaz Nadirov}

\textbf{Email:} a.nadirov@innopolis.university

\textbf{Area: Infrastructure and Deployment}

Responsibilities:

\begin{itemize}
    \item GitHub repository organization
    \item dependency and environment management
    \item Docker configuration
    \item GitHub Actions configuration
    \item automated testing
    \item deployment of the final application
    \item reproducible execution of experiments
    \item preparation of setup, training, evaluation, and launch instructions
\end{itemize}

\subsection{Marat Yunusov}

\textbf{Email:} mar.yunusov@innopolis.university

\textbf{Area: Frontend and Qualitative Analysis}

Responsibilities:

\begin{itemize}
    \item implementation of the React frontend
    \item communication between the frontend and backend
    \item image upload and prediction display
    \item interface states and error handling
    \item integration testing of the complete application
    \item preparation of result visualizations
    \item manual inspection of classification errors
    \item grouping representative failures by their likely visual cause
    \item preparation of representative prediction examples for the final demonstration
\end{itemize}

\end{document}